\documentclass[11pt]{article}
\usepackage[a4paper,margin=25mm]{geometry}
\usepackage[T1]{fontenc}
\usepackage{lmodern,amsmath,amssymb,amsthm,mathtools,microtype}
\usepackage[colorlinks=true,linkcolor=blue,citecolor=blue,urlcolor=blue]{hyperref}
\newtheorem{proposition}{Proposition}[section]
\newtheorem{corollary}[proposition]{Corollary}
\newtheorem{lemma}[proposition]{Lemma}
\theoremstyle{remark}
\newtheorem{remark}[proposition]{Remark}
\newcommand{\Ric}{\operatorname{Ric}}
\newcommand{\Om}{\Omega}
\newcommand{\dd}{\,dV_g}
\title{Q4: Canonical Orthogonality and a Positive-Multiplier Criterion}
\author{}
\date{12 October 2026}
\begin{document}
\maketitle
\begin{abstract}
We combine the canonical-divisor theorem for closed real
\(J\)-anti-invariant forms with the canonical curvature formula for
Einstein almost-K\"ahler four-manifolds. The transverse Weyl form is
\(L^2\)-orthogonal to every closed \(J\)-anti-invariant form. Consequently,
a globally smooth positive integrating factor for the transverse Weyl form
forces K\"ahlerness. No Betti-number or Seiberg--Witten chamber hypothesis
is needed. We also compute the obstruction to such a factor on the
project's explicit noncompact model: its Lee form has nonzero exterior
derivative on every regular open set. The general existence of a positive
factor on a closed Einstein candidate is not established.
\end{abstract}

\section{Data and the exact global pairing}
Let \((M^4,g,J,\Om)\) be closed and connected, with
\[
 d\Om=0,\qquad \Om(X,Y)=g(JX,Y),\qquad |\Om|^2=2,\qquad
 \Ric=\frac{s}{4}g.
\]
The orientation is that of \(J\), and \(s\) is constant.
Distinguish the Hodge splitting \(\Lambda_g^\pm\) from the
\(J\)-type splitting \(\Lambda_J^\pm\). In dimension four,
\[
 \Lambda_J^-\subset\Lambda_g^+,\qquad
 \Lambda_g^+=\mathbb R\Om\oplus\Lambda_J^-,
 \qquad
 \Lambda_J^+=\mathbb R\Om\oplus\Lambda_g^-.
\]
Write \(K=-c_1(TM,J)\), \(q=|\nabla\Om|^2\), and
\[
 \eta=W^+(\Om)-\left(\frac{s}{6}+\frac q4\right)\Om
 \in\Gamma(\Lambda_J^-).
\]
LeBrun's Proposition 2.1 \cite{LeBrun} gives the real canonical
Chern form \(\rho=iF_{\mathcal A}\), with
\[
 [\rho]=-2\pi K,\qquad
 \rho=\eta+\left(\frac s4+\frac q8\right)\Om-\frac q8\widehat\Om .
 \tag{1}
\]
Here \(\widehat\Om\) is anti-self-dual where \(q\ne0\).
The product \(q\widehat\Om\) extends smoothly through \(q=0\);
all subsequent expressions use that smooth product, not division by \(q\).
The project supplied exactly this normalization.

Set
\[
 \beta=\rho-\frac{s}{4}\Om,\qquad
 [\beta]=-2\pi K-\frac{s}{4}[\Om].
 \tag{2}
\]
Then \(d\beta=0\), and its \(J\)-anti-invariant part is \(\eta\).
The other part is \(J\)-invariant, so for every
\(\alpha\in\Gamma(\Lambda_J^-)\),
\[
 \beta\wedge\alpha=\langle\eta,\alpha\rangle\dd.
 \tag{3}
\]

\begin{lemma}[Canonical pairing]
On a closed connected almost-complex four-manifold,
every closed real \(J\)-anti-invariant form \(\alpha\) satisfies
\[
 K\cdot[\alpha]=0.
 \tag{4}
\]
\end{lemma}
\begin{proof}
There is nothing to prove if \(\alpha=0\).
If \(\alpha\) is nowhere zero, it trivializes the complex line underlying
\(\Lambda_J^-\), and hence \(K=0\).
Otherwise the canonical-divisor theorem of Bonthrone--Zhang
\cite[Theorem 1.1]{BZ} gives a finite \(J\)-holomorphic
one-subvariety \(\Theta_\alpha\), supported on the zero set of \(\alpha\),
whose Poincar\'e dual is \(K\). Thus
\[
 K\cdot[\alpha]=\int_{\Theta_\alpha}\alpha=0.
\]
Multiplicities and singular points do not change this evaluation.
\end{proof}

\begin{proposition}[Orthogonality of the transverse Weyl form]
For the closed Einstein almost-K\"ahler data above,
\[
 \boxed{\int_M\langle\eta,\alpha\rangle\dd=0
 \quad\text{for every }\alpha\in\Gamma(\Lambda_J^-),\ d\alpha=0.}
 \tag{5}
\]
\end{proposition}
\begin{proof}
Because \(\Om\wedge\alpha=0\) pointwise, equations (2)--(4) give
\[
 \int_M\langle\eta,\alpha\rangle\dd
 =[\beta]\cdot[\alpha]
 =-2\pi K\cdot[\alpha]-\frac{s}{4}[\Om]\cdot[\alpha]=0.
\]
The argument is independent of the value of \(b_2^+\).
\end{proof}

\section{A positive-multiplier criterion}
\begin{corollary}
Suppose \(f\in C^\infty(M)\), \(f>0\), and \(d(f\eta)=0\).
Then \(\eta=0\), and \((g,J,\Om)\) is K\"ahler.
In particular, \(d\eta=0\) is sufficient.
\end{corollary}
\begin{proof}
Use \(\alpha=f\eta\) in (5). It gives
\[
 0=\int_M f|\eta|^2\dd.
 \tag{6}
\]
Since \(f>0\), one has \(\eta=0\).
Thus \(\Om\) is an eigenform of \(W^+\).
The Einstein condition supplies \(J\)-invariant Ricci, so the compact
third-Gray theorem of Apostolov--Armstrong--Dr\u aghici
\cite[Theorem 2]{AAD} gives \(\nabla\Om=0\).
\end{proof}

This is a combination of the cited theorems with (1), not a proof that
the multiplier exists under the unaugmented negative Goldberg hypotheses.
It yields the exact equivalence, on the stated closed Einstein class,
\[
 \nabla\Om=0
 \ \Longleftrightarrow\ \eta=0
 \ \Longleftrightarrow\ d\eta=0
 \ \Longleftrightarrow\
 \exists f\in C^\infty(M),\ f>0,\ d(f\eta)=0.
 \tag{7}
\]
Existence of \(f\), not the implications in (7), remains the issue
when this criterion is used to approach Q4.

\section{What the multiplier equation really requires}
On \(U=\{|\eta|>0\}\), the form \(\eta\) is nondegenerate because
\(\eta\wedge\eta=|\eta|^2dV_g\).
There is a unique smooth one-form \(\vartheta_\eta\) on \(U\) such that
\[
 d\eta=\vartheta_\eta\wedge\eta.
 \tag{8}
\]
The map \(\xi\mapsto\xi\wedge\eta\) from one-forms to three-forms
is an isomorphism at each point of \(U\).
Therefore
\[
 d(f\eta)=0\quad\Longleftrightarrow\quad
 d\log f=-\vartheta_\eta\quad\text{on }U.
 \tag{9}
\]
Local existence requires \(d\vartheta_\eta=0\); on a non-simply-connected
component it also requires vanishing periods. For the corollary,
the resulting \(f\) must extend smoothly and positively over \(M\).
Although \(d^2\eta=0\) gives
\[
 d\vartheta_\eta\wedge\eta=0,
\]
this is only a primitive-curvature condition and does not say
\(d\vartheta_\eta=0\).
Likewise, if \(\sigma=\beta-\eta\), then
\[
 d\eta=-d\sigma
\]
follows from \(d\beta=0\), not \(d\eta=0\).

\section{Exact computation on the supplied noncompact model}
This section uses the project's existing metric and curvature blocks
\cite{Project,ProjectCurv}. It does not construct another example or
re-prove their Einstein equations.

Put \(1/3<x<1\), \(D=1+3x^2\), and use invariant one-forms with
\[
 d\sigma_1=-2\sigma_2\wedge\sigma_3,\quad
 d\sigma_2=2\sigma_3\wedge\sigma_1,\quad
 d\sigma_3=2\sigma_1\wedge\sigma_2.
\]
The source model has
\[
 \Om=\sqrt{X_1}\,\omega_1+\sqrt{X_2}\,\omega_2,\quad
 X_1=\frac{(1-x)^3(3x+1)}{16x^3},\quad
 X_2=\frac{(1+x)^3(3x-1)}{16x^3},
\]
where \(\omega_i=e^{0i}+e^{jk}\), and the source curvature eigenvalues satisfy
\[
 r_1-r_2=\frac{48x^3}{D^2}.
\]
Consequently
\[
 \eta=(r_1-r_2)
 \left(X_2\sqrt{X_1}\,\omega_1-X_1\sqrt{X_2}\,\omega_2\right).
 \tag{10}
\]
In the invariant coframe, set \(F=\sqrt3/(x^{3/2}D)\). Then
\[
 \eta=F\bigl(a\,dx\wedge\sigma_1+b\,dx\wedge\sigma_2
              +c\,\sigma_2\wedge\sigma_3+d\,\sigma_3\wedge\sigma_1\bigr),
 \tag{11}
\]
with
\[
 \begin{split}
 a&=(x+1)(3x-1),\qquad b=(x-1)(3x+1),\\
 c&=\frac{4x(x+1)^2(9x^2-1)}D,\qquad
 d=-\frac{4x(x-1)^2(9x^2-1)}D.
 \end{split}
\]
The form is nonzero throughout this regular interval; indeed
\[
 |\eta|^2=\frac{18(1-x^2)^3(9x^2-1)}{D^4}>0.
\]
A direct exterior differentiation of (11) gives
\[
 \boxed{
 \vartheta_\eta=
 -\frac{3(x^2-1)}{xD}\,dx+\frac{12(x^2-1)}D\,\sigma_3,
 }\tag{12}
\]
and hence
\[
 \boxed{
 d\vartheta_\eta=
 \frac{96x}{D^2}\,dx\wedge\sigma_3+
 \frac{24(x^2-1)}D\,\sigma_1\wedge\sigma_2\ne0.
 }\tag{13}
\]
For a reproducible verification, with
\(u=-3(x^2-1)/(xD)\), \(v=12(x^2-1)/D\) and
\(\ell=(\log F)'=-3/(2x)-6x/D\), the two nonzero coefficients of
\(d\eta-\vartheta_\eta\wedge\eta\), after division by \(F\), are
\[
 c'+\ell c+2a-uc-vb=0,\qquad
 d'+\ell d-2b-ud+va=0.
\]
These are exact rational identities. The remaining coefficients vanish.
Differentiating \(v\) gives \(v'=96x/D^2\), proving (13).

It follows that no nonvanishing smooth multiplier, radial or otherwise,
can make the source \(\eta\) closed on any nonempty regular open set.
This does not contradict the compact criterion: the source is noncompact.
It shows that a multiplier cannot follow from the local Einstein and
almost-K\"ahler equations alone. A proposed global proof must supply
additional information using closedness.

\section{The remaining Q4 connection}
The canonical-divisor input upgrades a nonnegative curve pairing to the
exact vanishing (5). It finishes the implication from a positive
closed multiple of \(\eta\) to K\"ahlerness, uniformly in the Betti-number
cases and without the project's FD or N1 estimates.
What is not supplied is a proof that an arbitrary closed negative
Einstein almost-K\"ahler candidate has such a multiplier, or an independent
argument forcing \([\beta]^2=0\).
The general closed negative Goldberg conclusion and the other three
original questions are not resolved by this note.

\begin{thebibliography}{9}
\bibitem{BZ}
L.~Bonthrone and W.~Zhang,
\emph{\(J\)-holomorphic curves from closed \(J\)-anti-invariant forms},
arXiv:1808.09356v2 (2020), Theorem 1.1.
\url{https://arxiv.org/abs/1808.09356}
\bibitem{LeBrun}
C.~LeBrun, \emph{Einstein Metrics, Symplectic Minimality, and
Pseudo-Holomorphic Curves}, arXiv:math/0409290, Proposition 2.1.
\url{https://arxiv.org/abs/math/0409290}
\bibitem{AAD}
V.~Apostolov, J.~Armstrong and T.~Dr\u aghici,
\emph{Local models and integrability of certain almost K\"ahler
four-manifolds}, Math.\ Ann.\ 323 (2002), 633--666, Theorem 2.
\url{https://arxiv.org/abs/math/0109098}
\bibitem{Project}
Supplied project manuscript,
\emph{03\_Subsequent\_Mathematical\_Developments(1).tex},
8 October 2026, Section 2: source metric, almost-K\"ahler form and invariant
coframe. The \(\beta\) normalization is also in
\emph{01\_Research\_Synthesis(1).tex}, Goldberg section.
\bibitem{ProjectCurv}
Supplied continuation,
\emph{Pair\_Gap\_and\_Pinched\_Example.tex},
8 October 2026, source curvature blocks.
\end{thebibliography}
\end{document}
